%==============
\begin{center}
{\LARGE \bf Quantifier Elimination
}
\end{center}
%==============
\bigskip

\noindent
{\bf Organizer}
\medskip

\begin{itemize}

\item
S. Steinberg (stanly@math.unm.edu)
\item
H. Hong (Hoon.Hong@risc.uni-linz.ac.at)

\item
R. Liska (liska@siduri.fjfi.cvut.cz)

\end{itemize}

\bigskip

\noindent
{\bf Description}
\medskip

Recently, quantifier elimination algorithms have seen many applications, and there have been new
development in quantifier elimination algorithms for special classes of problems. 

\bigskip
\noindent
{\bf Talks}
\medskip

Date: July 19th (Friday) 

\begin{description}
\item[14:00-14:20] \ \\
  Andreas Dolzmann, Thomas Sturm\\
  {\bf REDLOG -- Computer Algebra Meets Computer Logic.} 
\item[14:20-14:40] \ \\
  Volker Weispfenning\\
  {\bf Quantifier Elimination Applied in Simulation and Scheduling} 
\item[14:40-15:00] \ \\
  Emma-Neila Gonzalez-Campos, Laureano Gonzalez-Vega\\
  {\bf Global projection operators for the Cylindrical Algebraic Decomposition algorithm.} 
\item[15:00-15:20] \ \\
  Chris Brown, George Collins\\
  {\bf Simple Solution Formula Construction from Truth Invariant CADs.} 
\item[15:20-15:50] \ \\
  \\
  {\bf Coffee Break} 
\item[15:50-16:10] \ \\
  Mats Jirstrand\\
  {\bf Some Applications of Quantifier Elimination in Nonlinear
    Control Theory.} 
\item[16:10-16:30] \ \\
  Ying Jiangqian\\
  {\bf On the Determination of the Strong Stabilizability of an
    n-Dimensional Linear System.} 
\item[16:30-16:50] \ \\
  H. Hong, A. Neubacher\\
  {\bf Approximate Quantifier Elimination.} 
\item[16:50-17:10] \ \\
  R. Liska, S. Steinberg\\
  {\bf The Stability of Boundary Conditions.} 
\item[17:10-17:30] \ \\
  H. Hong, R. Liska, S. Steinberg\\
  {\bf Discussion and Comments by the Organizers and Participants. } 
\end{description}

